\section*{CHƯƠNG 1: CƠ SỞ LÝ LUẬN VÀ BỐI CẢNH LỊCH SỬ -- XÃ HỘI TÁC ĐỘNG ĐẾN TẦNG LỚP NHO SĨ VIỆT NAM THẾ KỶ XVIII}
\addcontentsline{toc}{section}{CHƯƠNG 1: CƠ SỞ LÝ LUẬN VÀ BỐI CẢNH LỊCH SỬ -- XÃ HỘI TÁC ĐỘNG ĐẾN TẦNG LỚP NHO SĨ VIỆT NAM THẾ KỶ XVIII}

\subsection*{1.1. Quan điểm triết học Mác -- Lênin về ý thức xã hội và khái lược diện mạo nho sĩ trước thế kỷ XVIII}
\addcontentsline{toc}{subsection}{1.1. Quan điểm triết học Mác -- Lênin về ý thức xã hội và khái lược diện mạo nho sĩ trước thế kỷ XVIII}

\subsubsection*{1.1.1. Nguyên lý duy vật lịch sử về tồn tại xã hội, ý thức xã hội và đặc điểm tầng lớp trí thức}
\addcontentsline{toc}{subsubsection}{1.1.1. Nguyên lý duy vật lịch sử về tồn tại xã hội, ý thức xã hội và đặc điểm tầng lớp trí thức}

Chủ nghĩa duy vật lịch sử giải thích sự vận động của xã hội loài người xuất phát từ mối quan hệ biện chứng giữa tồn tại xã hội và ý thức xã hội. Tồn tại xã hội là toàn bộ phương diện sinh hoạt vật chất và các điều kiện sinh hoạt vật chất của xã hội, bao gồm phương thức sản xuất vật chất, điều kiện tự nhiên -- địa lý và dân cư, trong đó phương thức sản xuất giữ vai trò quyết định. Ý thức xã hội là phương diện sinh hoạt tinh thần của xã hội, nảy sinh từ tồn tại xã hội và phản ánh tồn tại xã hội trong những giai đoạn lịch sử nhất định thông qua các hình thái như ý thức chính trị, pháp quyền, đạo đức, tôn giáo, nghệ thuật, khoa học và triết học.

Các Mác và Phriđrích Ăngghen khẳng định tồn tại xã hội quyết định nội dung, tính chất và xu hướng biến đổi của ý thức xã hội; khi phương thức sản xuất và các quan hệ kinh tế thay đổi thì những tư tưởng, quan điểm chính trị, đạo đức sớm muộn cũng biến đổi theo. Song học thuyết mácxít đồng thời chỉ rõ ý thức xã hội có tính độc lập tương đối, thể hiện trên bốn khía cạnh: ý thức xã hội thường lạc hậu hơn tồn tại xã hội do sức ỳ của tập quán và sự níu kéo từ lợi ích của các giai cấp thống trị cũ; một số tư tưởng khoa học, tiến bộ có thể vượt trước, dự báo xu hướng vận động của hiện thực; sự phát triển của ý thức xã hội luôn mang tính kế thừa lịch sử; và đặc biệt, các hình thái ý thức xã hội có sự tương tác lẫn nhau cũng như tác động trở lại tồn tại xã hội theo hai chiều hướng thúc đẩy hoặc kìm hãm, tùy thuộc vào mức độ phù hợp của hệ tư tưởng đó với quy luật khách quan.

Trong kết cấu xã hội -- giai cấp, triết học Mác -- Lênin xác định trí thức là một tầng lớp xã hội đặc thù chuyên lao động trí óc phức tạp, trực tiếp sáng tạo, bảo lưu và truyền bá tri thức văn hóa, tư tưởng. Trí thức không cấu thành một giai cấp kinh tế độc lập vì họ không nắm giữ tư liệu sản xuất riêng biệt và không đại diện cho một phương thức sản xuất riêng. Chính vị trí trung gian về mặt kinh tế khiến lập trường chính trị và thế giới quan của trí thức luôn chịu sự quy định bởi các giai cấp cơ bản trong xã hội. Khi chế độ xã hội ổn định và giai cấp cầm quyền còn đóng vai trò tiến bộ, trí thức là lực lượng nòng cốt củng cố thiết chế và hệ tư tưởng chính thống. Ngược lại, khi phương thức sản xuất cũ lâm vào bế tắc và mâu thuẫn giai cấp gay gắt, tầng lớp trí thức tất yếu mất đi tính thống nhất nội tại, buộc phải lựa chọn lại chỗ đứng chính trị và định hướng học thuật trước những đòi hỏi mới của lịch sử.

\subsubsection*{1.1.2. Khái niệm nho sĩ và vị thế của tầng lớp nho sĩ Việt Nam trong lịch sử trung đại trước thế kỷ XVIII}
\addcontentsline{toc}{subsubsection}{1.1.2. Khái niệm nho sĩ và vị thế của tầng lớp nho sĩ Việt Nam trong lịch sử trung đại trước thế kỷ XVIII}

Trong xã hội phong kiến Việt Nam, nho sĩ (hay kẻ sĩ, sĩ phu Nho học) là tầng lớp trí thức được đào tạo bằng hệ thống kinh điển Nho gia (\textit{Tứ thư}, \textit{Ngũ kinh}), tham gia khoa cử để bước vào bộ máy quan liêu hoặc lui về làng xã dạy học, bốc thuốc, giữ gìn hương ước. Trong trật tự phân tầng ``Sĩ, Nông, Công, Thương'', kẻ sĩ đứng ở vị trí đầu bảng nhờ uy tín học vấn và vai trò biểu trưng đạo lý. Nhân sinh quan chính trị của họ xoay quanh lộ trình ``tu thân, tề gia, trị quốc, bình thiên hạ'' và nguyên tắc ứng xử ``xuất -- xử'': gặp thời thịnh trị, vua sáng tôi hiền thì xuất thế hành đạo giúp nước; gặp buổi hôn quân bạo chúa, kỷ cương đảo lộn thì ẩn cư giữ trọn khí tiết người quân tử.

Nhìn lại chặng đường trước thế kỷ XVIII, vị thế của tầng lớp nho sĩ gắn liền với quá trình xác lập và biến đổi của Nho giáo tại Đại Việt. Dưới triều Lý -- Trần (thế kỷ XI -- XIV), Nho giáo từng bước đi vào thiết chế giáo dục và hành chính qua việc xây dựng Văn Miếu (1070), mở khoa thi Minh kinh bác học đầu tiên (1075) và lập Quốc Tử Giám (1076). Tuy vậy, trong suốt thời kỳ này, Nho giáo chưa giữ địa vị độc tôn mà dung hợp hài hòa với Phật giáo và Đạo giáo trên nền tảng ``Tam giáo đồng nguyên'', tạo nên lớp nho sĩ vừa mẫn cán việc triều chính vừa cởi mở, phóng khoáng về tư duy.

Sang thời Lê sơ thế kỷ XV, đỉnh cao là triều vua Lê Thánh Tông, nhà nước quân chủ trung ương tập quyền đưa Nho giáo (theo hệ phái Tống Nho của Trình -- Chu) lên địa vị độc tôn làm bệ đỡ tư tưởng chính thống. Bài ký trên bia Tiến sĩ khoa Nhâm Tuất dựng tại Văn Miếu năm 1484 do Thân Nhân Trung soạn đã đúc kết quan điểm cốt lõi của thời đại: ``Hiền tài là nguyên khí của quốc gia, nguyên khí thịnh thì thế nước mạnh rồi lên cao, nguyên khí suy thì thế nước yếu rồi xuống thấp''. Nho sĩ khoa bảng trở thành rường cột trực tiếp vận hành nhà nước pháp chế Hồng Đức, đại diện cho kỷ cương phong kiến đang ở giai đoạn hưng thịnh.

Đến thế kỷ XVI -- XVII, sự suy yếu của nhà Lê sơ đẩy đất nước vào cục diện phân tranh kéo dài giữa Nam -- Bắc triều (Lê -- Mạc) và tiếp đó là Trịnh -- Nguyễn phân tranh chia cắt hai miền Đàng Ngoài, Đàng Trong. Thực tế vua Lê chỉ còn hư vị, quyền bính thuộc về các chúa đã làm rạn nứt quan niệm ``chính danh'' và thuyết ``thiên mệnh'' Nho gia. Một số bậc đại nho tiêu biểu như Trạng Trình Nguyễn Bỉnh Khiêm đã chủ động cáo quan về dựng am Bạch Vân dạy học, bày tỏ thái độ phản tỉnh trước các cuộc nội chiến huynh đệ tương tàn. Dù vậy, trong hai thế kỷ này, hệ chuẩn luân lý tam cương ngũ thường về cơ bản vẫn kiểm soát được trật tự xã hội truyền thống.

\subsection*{1.2. Điều kiện kinh tế, chính trị, văn hóa thế kỷ XVIII và những biến động tác động đến tầng lớp nho sĩ}
\addcontentsline{toc}{subsection}{1.2. Điều kiện kinh tế, chính trị, văn hóa thế kỷ XVIII và những biến động tác động đến tầng lớp nho sĩ}

\subsubsection*{1.2.1. Sự khủng hoảng toàn diện của thể chế phong kiến Đàng Ngoài -- Đàng Trong và sự suy đồi của khoa cử}
\addcontentsline{toc}{subsubsection}{1.2.1. Sự khủng hoảng toàn diện của thể chế phong kiến Đàng Ngoài -- Đàng Trong và sự suy đồi của khoa cử}

Bước sang thế kỷ XVIII, chế độ phong kiến ở cả hai miền Nam -- Bắc đồng loạt lâm vào khủng hoảng cấu trúc. Tại Đàng Ngoài, thể chế lưỡng đầu Vua Lê -- Chúa Trịnh bộc lộ sự mục ruỗng từ thượng tầng: danh nghĩa ``phù Lê'' chỉ còn là vỏ bọc cho nạn chuyên quyền, kiêu binh nổi loạn và những cuộc thanh trừng phe phái trong phủ chúa. Ở địa phương, quan lại cấu kết với địa chủ cường hào ráo riết kiêm tính ruộng đất công tư, đẩy nông dân vào cảnh mất tư liệu sản xuất, bần cùng hóa và lưu tán hàng loạt giữa những nạn đói lớn. Tại Đàng Trong, sau giai đoạn đầu khai phá năng động, chính quyền chúa Nguyễn nửa sau thế kỷ XVIII cũng suy đồi nhanh chóng, điển hình dưới thời chúa Nguyễn Phúc Thuần khi quyền thần Trương Phúc Loan thao túng triều chính, mua quan bán tước và đặt ra chế độ thuế khóa nặng nề.

Trên phương diện kinh tế -- xã hội, bất chấp chính sách ``trọng nông ức thương'' của nhà nước phong kiến, kinh tế hàng hóa và luồng giao thương đô thị ở thế kỷ XVIII (tại Kẻ Chợ, Phố Hiến, Hội An, Cù lao Phố, Gia Định) vẫn có bước phát triển khởi sắc. Sự thâm nhập của quan hệ tiền tệ vào đời sống xã hội một mặt kích thích nhu cầu tri thức thực tiễn và ý thức thị dân, mặt khác khi kết hợp với bộ máy quan liêu suy đồi lại làm trầm trọng thêm tệ nạn thương mại hóa quyền lực và học vị.

Biểu hiện trực diện nhất tác động đến tâm thức nho sĩ là sự suy thoái của nền giáo dục và khoa cử quan phương. Quốc Tử Giám cùng hệ thống trường học phủ, huyện sa sút, kỷ cương trường thi buông lỏng với nạn mang sách, mướn người làm bài thay. Đỉnh điểm diễn ra vào năm Canh Ngọ 1750 dưới thời vua Lê Hiển Tông -- chúa Trịnh Doanh: do ngân khố kiệt quệ vì chiến phí đàn áp khởi nghĩa nông dân, triều đình chấp thuận đề xuất của Tham tụng Đỗ Thế Giai, cho phép người dự thi chỉ cần nộp ba quan ``tiền thông kinh'' là được miễn kỳ khảo hạch ở địa phương để vào thẳng trường thi Hương.

Việc xóa bỏ khâu sát hạch năng lực đầu vào khiến trường thi năm đó hỗn loạn chưa từng thấy; người buôn bán, kẻ giàu có không thông chữ nghĩa cũng ùn ùn nộp tiền vào trường chép mướn, dẫn đến cảnh đỗ đạt tràn lan và sinh ra câu châm biếm dân gian ``Sinh đồ ba quan''. Khi con đường khoa bảng -- vốn được kẻ sĩ tôn thờ là thước đo tài đức công minh -- bị đồng tiền thao túng công khai, niềm tin của giới nho sĩ chân chính vào tính chính đáng của triều đình và trật tự lễ giáo chính thống đã đổ vỡ nghiêm trọng.

\subsubsection*{1.2.2. Cao trào khởi nghĩa nông dân với đỉnh cao phong trào Tây Sơn và yêu cầu thời cuộc đối với kẻ sĩ}
\addcontentsline{toc}{subsubsection}{1.2.2. Cao trào khởi nghĩa nông dân với đỉnh cao phong trào Tây Sơn và yêu cầu thời cuộc đối với kẻ sĩ}

Áp bức giai cấp nặng nề đã châm ngòi cho các cuộc khởi nghĩa nông dân nổ ra liên tục suốt thế kỷ XVIII, làm lung lay tận gốc nền thống trị phong kiến. Ở Đàng Ngoài, các cuộc nổi dậy quy mô lớn của Nguyễn Danh Phương, Nguyễn Tuyển, Nguyễn Cừ, Hoàng Công Chất và đặc biệt là cuộc khởi nghĩa kéo dài hơn một thập niên của Nguyễn Hữu Cầu đã uy hiếp trực tiếp kinh thành Thăng Long, buộc giới sĩ phu phải chứng kiến sức mạnh phản kháng của tầng lớp nông dân vốn bị coi thường dưới đáy xã hội.

Đến năm 1771, cuộc khởi nghĩa Tây Sơn do ba anh em Nguyễn Nhạc, Nguyễn Lữ, Nguyễn Huệ phát động từ ấp Tây Sơn (Bình Định) đã chuyển hóa từ phong trào đấu tranh giai cấp ở địa phương thành cao trào thống nhất đất nước và bảo vệ độc lập dân tộc. Nghĩa quân Tây Sơn lần lượt lật đổ chính quyền chúa Nguyễn ở Đàng Trong (1777), đánh tan 5 vạn quân xâm lược Xiêm tại trận Rạch Gầm -- Xoài Mút (1785), rồi tiến ra Bắc Hà tiêu diệt thế lực chúa Trịnh (1786), chấm dứt hơn hai thế kỷ chia cắt sông Gianh. Đỉnh cao là mùa Xuân năm Kỷ Dậu 1789, dưới sự thống lĩnh của Hoàng đế Quang Trung -- Nguyễn Huệ, đại quân Tây Sơn đã hành quân thần tốc quét sạch 29 vạn quân Mãn Thanh do vua Lê Chiêu Thống rước về giày xéo đất nước, giữ vững nền độc lập tự chủ của Đại Việt.

Sự sụp đổ của các tập đoàn phong kiến Lê -- Trịnh, chúa Nguyễn cùng sự xác lập vương triều mới xuất thân từ áo vải nông dân đã tạo nên cú chấn động lớn đối với toàn bộ tầng lớp nho sĩ. Hệ chuẩn luân lý trung quân cứng nhắc của Tống Nho không còn chỗ dựa thực tế khi chính ông vua cuối cùng của triều Lê đã phản bội Tổ quốc, còn sứ mệnh cứu nước lại thuộc về những người anh hùng nông dân. Thực tại khách quan ấy trở thành phép thử lịch sử khắc nghiệt, thúc đẩy quá trình phân hóa chính trị và tái cấu trúc tư duy học thuật của nho sĩ Việt Nam trong suốt thế kỷ XVIII.
